% Part: methods
% Chapter: proofs
% Section: example-1

\documentclass[../../../include/open-logic-section]{subfiles}

\begin{document}

\olfileid{mth}{prf}{ex1}

\olsection{اک مثال}

ساڈی پہلی مثال سیٹاں دے اتحاد تے سانجھے حصے بارے ہِیٹھلی سوکھی گل اے۔ ایہ تعریفاں کھولھنا، عطف تے ساریاں لئی دعوے ثابت کرنا، تے صورتاں وکھ کر کے ثبوت دینا وکھاوے گی۔

\begin{prop}
کسے وی سیٹاں $A$، $B$ تے $C$ لئی، $A \cup (B \cap C) = (A \cup B)
\cap (A \cup C)$
\end{prop}

آؤ ایہنوں ثابت کریے!{}

\begin{proof}
اسیں وکھاؤنا چاہندے آں کہ کسے وی سیٹاں $A$، $B$ تے $C$ لئی، $A \cup (B \cap
C) = (A \cup B) \cap (A \cup C)$
\begin{quote}
پہلاں قضیے دے بیان وچ «$=$» دی تعریف کھولھدے آں۔ چیتے کرو کہ سیٹاں دا اکّو ہونا ثابت کرن دا مطلب ایہ وکھاؤنا اے کہ سیٹاں دے !!{element}s اکّو نیں۔ یعنی $A \cup
(B \cap C)$ دے سارے !!{element}s، $(A \cup B) \cap (A \cup C)$ دے !!{element}s وی نیں، تے اُلٹے پاسے وی۔ «اُلٹے پاسے وی» دا مطلب اے کہ $(A
\cup B) \cap (A \cup C)$ دا ہر !!{element}، $A \cup (B \cap
C)$ دا !!a{element} وی ہونا چاہیدا اے۔ سو تعریف کھولھ کے سانوں پتا لگدا اے کہ اسیں اک عطف ثابت کرنا اے۔ آؤ ایہ لکھ لئیے:
\end{quote}
تعریف مطابق، $A \cup (B \cap C) = (A \cup B) \cap (A \cup C)$ اُدوں تے صرف اُدوں جدوں $A \cup (B \cap C)$ دا ہر !!{element}، $(A
\cup B) \cap (A \cup C)$ دا !!a{element} وی اے، تے $(A \cup B) \cap (A
\cup C)$ دا ہر !!{element}، $A \cup (B \cap C)$ دا !!a{element} اے۔
\begin{quote}
کیوں جے ایہ عطف اے، سانوں اوہدے دوویں جز وکھ وکھ ثابت کرنے پَین گے۔ پہلاں پہلا جز لئیے۔ آؤ ثابت کریے کہ $A \cup (B \cap C)$ دا ہر !!{element}، $(A
\cup B) \cap (A \cup C)$ دا !!a{element} وی اے۔

ایہ ساریاں لئی دعویٰ اے۔ سو اسیں $A \cup (B \cap C)$ دا اک من مانا !!{element} لے کے وکھاؤندے آں کہ اوہ $(A \cup B) \cap (A \cup C)$ دا !!a{element} وی ہونا چاہیدا اے۔ ایس من مانے !!{element} دا ناں رکھن لئی اک متغیر چُناں گے، منّ لو~$z$۔ ساڈا ثبوت اگے ایویں چلدا اے:
\end{quote}
پہلاں اسیں ثابت کردے آں کہ $A \cup (B \cap C)$ دا ہر !!{element}، $(A \cup B) \cap (A \cup C)$ دا !!a{element} وی اے۔ $z \in A \cup (B
\cap C)$ ہووے۔ سانوں وکھاؤنا اے کہ $z \in (A \cup B) \cap (A \cup C)$۔
\begin{quote}
ہُن $\cup$ تے~$\cap$ دی تعریف کھولھن دا ویلا اے۔ مثلاً $\cup$ دی تعریف ایہ اے، $A \cup B = \Setabs{z}{z \in A
  \text{ یا } z \in B}$۔ جدوں اسیں «$A \cup (B
\cap C)$» اُتے ایہ تعریف لاؤندے آں، تاں تعریف وچ «$B$» والی تھاں ہُن «$B \cap C$» اے، سو $A \cup (B \cap C) = \Setabs{z}{z \in A \text{ یا }
  z \in B \cap C}$۔ سو ساڈے منّن دا مطلب کہ $z \in A \cup (B \cap C)$، ایہ اے کہ $z \in \Setabs{z}{z \in A \text{ یا } z \in B \cap
  C}$۔ تے $z \in \Setabs{z}{\dots z\dots}$ اُدوں تے صرف اُدوں جدوں \dots $z$ \dots، یعنی ایس صورت وچ، $z \in A$ یا $z \in B \cap C$۔
\end{quote}
$\cup$ دی تعریف مطابق یا $z \in A$ یا $z \in B \cap C$۔
\begin{quote}
ایہ فصل اے، سو صورتاں وکھ کر کے ثبوت دینا کم آوے گا۔ اسیں دوویں صورتاں لیندے آں تے ہر اک وچ وکھاؤندے آں کہ ساڈا لوڑیندا نتیجہ، یعنی «$z \in (A \cup B) \cap (A \cup
C)$»، نکلدا اے۔
\end{quote}
صورت ۱۔ منّ لو $z \in A$۔
\begin{quote}
منّیاں شرطاں توں ہور بہت کجھ کرن نوں نہیں۔ سو ویکھیے کہ نتیجے وچ ساڈے کم دی کیہڑی گل اے۔ اسیں وکھاؤنا چاہندے آں کہ $z \in (A \cup B) \cap (A \cup C)$۔ $\cap$ دی تعریف مطابق جے اسیں $z \in (A \cup B) \cap (A \cup C)$ وکھاؤنا اے، تاں وکھاؤنا پئے گا کہ ایہ $(A \cup B)$ تے $(A \cup C)$ دوہاں وچ اے۔ پر $z
\in A \cup B$ اُدوں تے صرف اُدوں جدوں $z \in A$ یا $z \in B$، تے ساڈے کول پہلاں ای صورت~۱ دی شرط اے کہ $z \in A$۔ ایس دلیل وچ $C$ نوں $B$ دی تھاں رکھیے تاں $z \in A \cup C$۔ ایہ دلیل اُلٹی سمت وچ چلی اے، سو آؤ اینوں ثبوت لئی لوڑیندی سمت وچ لکھیے۔
\end{quote}
کیوں جے $z \in A$، سو $z \in A$ یا $z \in B$، تے ایس لئی تعریف~$\cup$ مطابق $z \in A \cup B$۔ ایویں ای $z \in A \cup C$۔ پر ایہدا مطلب اے کہ $z \in (A \cup B) \cap (A \cup C)$، تعریف~$\cap$ مطابق۔
\begin{quote}
ایس نال صورتاں وکھ کر کے دتے ثبوت دی پہلی صورت پوری ہوئی۔ ہُن اسیں دوجی صورت وچ نتیجہ کڈھنا چاہندے آں، جتھے $z \in B \cap C$۔
\end{quote}
صورت ۲۔ منّ لو $z \in B \cap C$۔
\begin{quote}
اسیں فیر دو سیٹاں دے سانجھے حصے نال کم کردے آں۔ آؤ~$\cap$ دی تعریف لائیے:
\end{quote}
کیوں جے $z \in B \cap C$، سو $z$، $B$ تے $C$ دوہاں دا !!a{element} ہونا چاہیدا اے، تعریف~$\cap$ مطابق۔
\begin{quote}
ہُن فیر اپنے نتیجے ول ویکھن دا ویلا اے۔ سانوں وکھاؤنا اے کہ $z$، $(A \cup B)$ تے $(A \cup C)$ دوہاں وچ اے۔ فیر حل سِدھا ای نکلدا اے۔
\end{quote}
کیوں جے $z \in B$، سو $z \in (A \cup B)$۔ کیوں جے $z \in C$، سو $z \in (A
\cup C)$ وی۔ ایس لئی $z \in (A \cup B) \cap (A \cup C)$۔
\begin{quote}
ایتھے اسیں $\cup$ تے $\cap$ دیاں تعریفاں فیر لائیاں، پر کیوں جے اسیں اوہ پہلاں چیتے کر چکے آں تے وکھا چکے آں کہ جے $z$ دو سیٹاں وچوں کسے اک وچ اے تاں اوہ اوہناں دے اتحاد وچ اے، ایتھے ہر قدم کھولھ کے دسّن دی لوڑ نہیں۔

صورتاں وکھ کر کے ثبوت دی دوجی صورت وی پوری ہو گئی، سو ہُن اسیں اپنا پہلا نتیجہ دَس سکدے آں۔
\end{quote}
سو جے $z \in A \cup (B \cap C)$ تاں $z \in (A \cup B) \cap (A \cup C)$۔
\begin{quote}
ہُن بس دوجی سمت وکھاؤنی اے کہ $(A \cup B) \cap (A \cup C)$ دا ہر !!{element}، $A \cup (B \cap
C)$ دا !!a{element} اے۔ پہلاں وانگ ای اسیں پہلے سیٹ دا اک من مانا !!{element} منّ کے ساریاں لئی ایہ دعویٰ ثابت کردے آں، تے وکھاؤندے آں کہ اوہ دوجے سیٹ وچ وی ہونا چاہیدا اے۔ آؤ ایہ گل لکھ لئیے۔
\end{quote}
ہُن منّ لو $z \in (A \cup B) \cap (A \cup C)$۔ اسیں وکھاؤنا چاہندے آں کہ $z \in A \cup (B \cap C)$۔
\begin{quote}
ہُن اسیں ایس مفروضہ شرط توں کم کردے آں کہ $z \in (A \cup B) \cap (A
\cup C)$۔ آس اے کہ پہلے حصے والا اوہو~$z$ ورتن نال بھلیکھا نہیں پئے گا۔ جدوں اوہ حصہ مکیا سی، اوتھے منّیاں ساریاں شرطاں اوتھے ای مُک گئیاں۔ سو ہُن اسیں $z$ بارے نویاں شرطاں منّ سکدے آں۔ جے ایہ نال بھلیکھا پئے، تاں اگے $z$ دی تھاں کوئی ہور متغیر رکھ لو۔

سانوں پتا اے کہ $z$، $A \cup B$ تے $A \cup C$ دوہاں وچ اے، تعریف~$\cap$ مطابق۔ تے $\cup$ دی تعریف نال ایہنوں ہور کھولھ سکدے آں۔ یا $z \in A$ یا $z \in B$، تے نال یا $z \in A$ یا $z \in
C$۔ ایہ فیر صورتاں وکھ کر کے ثبوت لگدا اے، پر «تے» نال بھلیکھا پے سکدا اے۔ تسی سوچ سکدے او کہ ایہدا مطلب تِن صورتاں نیں، $z$ یا $A$، یا $B$، یا $C$ وچ اے۔ پر ایہ غلط ہووے گا۔ سانوں دھیان نال ہر فصل واری واری ویکھنی چاہیدی اے۔
\end{quote}
$\cap$ دی تعریف مطابق $z \in A \cup B$ تے $z \in A \cup C$۔ $\cup$ دی تعریف مطابق $z \in A$ یا $z \in B$۔ اسیں صورتاں وکھ کردے آں۔
\begin{quote}
کیوں جے اسیں پہلی فصل ول دھیان دیندے آں، اسیں (تعریف $A \cup C$ کھولھ کے) دوجی فصل ہن تک نہیں کڈھی۔ اصل وچ ہن تک اوہدی لوڑ وی نہیں۔ پہلی صورت $z \in A$ اے، تے کسے سیٹ دا !!a{element} اوس سیٹ دے کسے ہور سیٹ نال اتحاد دا !!a{element} وی اے۔ سو صورت~۱ سوکھی اے:
\end{quote}
صورت ۱۔ منّ لو $z \in A$۔ ایس توں نکلدا اے کہ $z \in A \cup (B \cap
C)$۔
\begin{quote}
ہُن دوجی صورت، $z \in B$۔ ایتھے دوجے $\cup$ دی تعریف کھولھاں گے تے فیر صورتاں وکھ کر کے ثبوت دیواں گے:
\end{quote}
صورت ۲۔ منّ لو $z \in B$۔ کیوں جے $z \in A \cup C$، سو یا $z \in
A$ یا $z \in C$۔ اسیں اگوں صورتاں وکھ کردے آں:

صورت ۲ا۔ $z \in A$۔ تاں فیر $z \in A \cup (B \cap C)$۔
\begin{quote}
ایہ ذرا اوپری گل سی۔ ایس صورت لئی سانوں ایہ منّن دی لوڑ اصل وچ نہیں پئی کہ~$z \in B$، پر ایس نال کوئی اوکڑ نہیں۔
\end{quote}
صورت ۲ب۔ $z \in C$۔ تاں $z \in B$ تے $z \in C$، سو $z \in B \cap C$، تے ایس لئی $z \in A \cup (B \cap C)$۔
\begin{quote}
ایس نال صورتاں وکھ کر کے دتے دوویں ثبوت پورے ہو گئے، سو دوجا ادھ وی پورا ہو گیا۔
\end{quote}
سو جے $z \in (A \cup B) \cap (A \cup C)$ تاں $z \in A \cup (B \cap C)$۔
\end{proof}

\end{document}
